% concatenated from IEEE-conference-template-062824.tex
\documentclass[conference]{IEEEtran}
\IEEEoverridecommandlockouts
% The preceding line is only needed to identify funding in the first footnote. If that is unneeded, please comment it out.
%Template version as of 6/27/2024

\usepackage{cite}
\usepackage{amsmath,amssymb,amsfonts}
\usepackage{amsthm}
\usepackage{algorithmic}
\usepackage{graphicx}
\usepackage{textcomp}
\usepackage{xcolor}

\usepackage{hyperref}
\usepackage{booktabs}

% Define theorem environments
\newtheorem{theorem}{Theorem}[section]
\newtheorem{lemma}[theorem]{Lemma}
\newtheorem{corollary}[theorem]{Corollary}


\def\BibTeX{{\rm B\kern-.05em{\sc i\kern-.025em b}\kern-.08em
    T\kern-.1667em\lower.7ex\hbox{E}\kern-.125emX}}
\begin{document}

\title{Spectral Analysis of the \\ Weighted Frobenius Objective}%*\\
% {\footnotesize \textsuperscript{*}Note: Sub-titles are not captured for https://ieeexplore.ieee.org  and
% should not be used}
% \thanks{Identify applicable funding agency here. If none, delete this.}
% }

\author{\IEEEauthorblockN{Vladislav Trifonov}
\IEEEauthorblockA{\textit{AI4S Center, Sberbank of Russia} \\
\textit{CAI, Skoltech}\\
Moscow, Russian Federation \\
vladislav.trifonov@skoltech.ru}
\and
% \IEEEauthorblockN{Oleg Iliev}
% % \IEEEauthorblockA{\textit{Flow and Material Simulation dept.} \\
% \IEEEauthorblockA{\textit{Fraunhofer ITWM} \\
% % \textit{Fraunhofer Institute for Industrial Mathematics ITWM}\\
% Kaiserslautern, Germany \\
% oleg.iliev@itwm.fraunhofer.de}
% \and
\IEEEauthorblockN{Ivan Oseledets}
\IEEEauthorblockA{\textit{AIRI} \\
\textit{CAI, Skoltech} \\
% \textit{name of organization (of Aff.)}\\
Moscow, Russia \\
oseledets@airi.net}
\and
\IEEEauthorblockN{Ekaterina Muravleva}
\IEEEauthorblockA{\textit{AI4S Center, Sberbank of Russia} \\
\textit{CAI, Skoltech}\\
Moscow, Russian Federation \\
e.muravleva@skoltech.ru}
}


\maketitle

\begin{abstract}
    We analyze a weighted Frobenius loss for approximating symmetric positive definite matrices in the context of preconditioning iterative solvers. Unlike the standard Frobenius norm, the weighted loss penalizes error components associated with small eigenvalues of the system matrix more strongly. Our analysis reveals that each eigenmode is scaled by the corresponding square of its eigenvalue, and that, under a fixed error budget, the loss is minimized only when the error is confined to the direction of the largest eigenvalue. This provides a rigorous explanation of why minimizing the weighted loss naturally suppresses low-frequency components, which can be a desirable strategy for the conjugate gradient method. The analysis is independent of the specific approximation scheme or sparsity pattern, and applies equally to incomplete factorizations, algebraic updates, and learning-based constructions. Numerical experiments confirm the predictions of the theory, including an illustration where sparse factors are trained by a direct gradient updates to IC(0) factor entries, i.e., no trained neural network model is used.
\end{abstract}

\begin{IEEEkeywords}
    Preconditioning; Conjugate Gradient Method; Sparse Matrix Approximations; Spectral Properties
\end{IEEEkeywords}


\section{Motivation and background}

We are concerned with the construction of structured approximations to large sparse matrices in the context of preconditioning. Preconditioners of this type are widely used to accelerate iterative methods for solving linear systems, especially those arising from discretizations of partial differential equations. While our focus is on preconditioners, the considerations developed here may also be relevant for other approximation problems where a structured approximation of a matrix is required.

A common procedure is to select a structured matrix $P$ that minimizes the Frobenius functional~[1]
\begin{equation}
    \label{eq:unweighted-frobenius}
\|P - A\|_F^2,
\end{equation}
with $A$ the given matrix. This classical choice is orthogonally invariant and computationally convenient. At the same time, this functional assigns equal importance to all components of the error, without regard to the spectral properties of $A$.

In this paper we investigate the weighted Frobenius functional
\begin{equation}
    \label{eq:weighted-frobenius}
\|(P - A)A^{-1}\|_F^2,
\end{equation}
where $A$ is assumed symmetric positive definite and $P$ is again restricted to a structured class, e.g., sparse factors with a prescribed sparsity pattern of incomplete Cholesky factorizations with zero level fill-in (IC(0)). Since $A$ is SPD, the eigendecomposition $A = Q \Lambda Q^\top$ exists. Expressing the error in this eigenbasis allows a precise analysis of how the two functionals differ.

The assumption that $A$ is symmetric positive definite is not restrictive in many applications. Such matrices arise naturally as stiffness matrices from finite element or finite difference discretizations of elliptic partial differential equations, as covariance matrices in statistics, and as normal equations in least-squares problems. In the PDE setting, the smallest eigenvalues are typically associated with smooth, slowly varying error modes, while the largest eigenvalues correspond to oscillatory components introduced by the discretization grid. The distinction between these modes is central in numerical analysis, as different iterative methods and preconditioners affect them in different ways.

Approximations $P$ with prescribed sparsity or structure are frequently used as preconditioners to accelerate iterative solvers for these systems. Classical examples include IC(0), which maintain the sparsity of the original matrix, and algebraic approximations designed to capture the dominant spectral behavior of $A$ while remaining computationally efficient. In such contexts, the choice of functional used to determine $P$ directly influences which parts of the spectrum are emphasized by the approximation.

\section{Spectral properties of Frobenius functionals}

We now analyze the difference between the unweighted functional~\eqref{eq:unweighted-frobenius}
and the weighted version~\eqref{eq:weighted-frobenius}, where $A$ is symmetric positive definite and $P$ is a structured approximation. The following lemma provides a spectral decomposition of both functionals.

\begin{lemma}[Spectral decomposition]\label{lem:spectral-objectives}
Let $A = Q \Lambda Q^\top$ be the eigen-decomposition of a symmetric positive definite matrix 
$A \in \mathbb{R}^{n \times n}$.%, with 
% $\Lambda = \mathrm{diag}(\lambda_1,\dots,\lambda_n)$ and 
% $0 < \lambda_1 \le \cdots \le \lambda_n$. 
For any matrix $P$, set $E := P - A$ and define
\begin{equation*}
a_j := \sum_{i=1}^n \big|(Q^\top E Q)_{ij}\big|^2, \qquad j=1,\dots,n.
\end{equation*}
Then
\begin{equation*}
\|E\|_F^2 = \sum_{j=1}^n a_j,
\qquad
\|E A^{-1}\|_F^2 = \sum_{j=1}^n \frac{a_j}{\lambda_j^2}.
\end{equation*}
\end{lemma}

\begin{proof}
Since $A^{-1} = Q \Lambda^{-1} Q^\top$, we have
\begin{equation*}
E A^{-1} = Q (Q^\top E Q) \Lambda^{-1} Q^\top.
\end{equation*}
By invariance of the Frobenius norm under orthogonal transformations,
\begin{equation*}
\|E\|_F^2 = \|Q^\top E Q\|_F^2 = \sum_{j=1}^n \sum_{i=1}^n |(Q^\top E Q)_{ij}|^2 = \sum_{j=1}^n a_j
\end{equation*}
and
\begin{align*}
\|E A^{-1}\|_F^2 &= \|(Q^\top E Q)\Lambda^{-1}\|_F^2 \nonumber \\
&= \sum_{j=1}^n \sum_{i=1}^n \frac{|(Q^\top E Q)_{ij}|^2}{\lambda_j^2} = \sum_{j=1}^n \frac{a_j}{\lambda_j^2}
\end{align*}
\end{proof}

\begin{theorem}[Weighted functional inequality]
For $A$, $E$, and $a_j$ as in Lemma~\ref{lem:spectral-objectives}, assume that 
$\sum_{j=1}^n a_j = c > 0$. Then
\begin{equation*}
\|E A^{-1}\|_F^2 = \sum_{j=1}^n \frac{a_j}{\lambda_j^2} \;\;\ge\;\; \frac{c}{\lambda_n^2},
\end{equation*}
with equality if and only if $a_j = 0$ for all $j < n$ and $a_n = c$.
\end{theorem}

\begin{proof}
Because $\lambda_1^{-2} \ge \cdots \ge \lambda_n^{-2}$ and the mapping 
$(a_1,\dots,a_n) \mapsto \sum_j a_j/\lambda_j^2$ is linear, 
the minimum under the constraint $\sum_j a_j = c$ is achieved by placing all energy 
at the smallest weight $\lambda_n^{-2}$. This proves the inequality and the 
corresponding condition for equality.
\end{proof}


\section{Computational experiments}

Recent advances in machine learning have explored neural network-based approaches for preconditioning, particularly Graph Neural Networks (GNNs) for constructing sparse factorized preconditioners~[3, 5]. However, recent studies~[4] have demonstrated fundamental limitations of message-passing GNNs in approximating non-local sparse triangular factorizations.

Our goal is to explore whether preconditioners can be improved by directly optimizing their sparse factors with gradient descent on different objectives. We investigate if this parameter-free approach can lead to better spectral properties for IC(0) factors with classical algorithm.

Our experiments aimed to demonstrate the effect on the spectral properties of the preconditioner for a single matrix. While this example is limited, many real-world applications require solving linear systems with a single coefficient matrix $A$ but multiple right-hand sides $\{ b_i\}_{i=1}^N$. This scenario is common in:

\begin{itemize}
\item Time-dependent PDEs, where the system matrix remains constant while the right-hand side evolves over time.
\item Parameter studies in optimization, where parameters change but the system structure remains fixed.
\item Multiple load cases in structural analysis, where different loading conditions are applied to the same structure.
\item Monte Carlo simulations, where random variations affect the right-hand side but not the system matrix.
\end{itemize}

Moreover, this approach operates globally on the entire preconditioner structure, which does not suffer from the principal limitations of GNNs for approximating sparse triangular factorizations.

To be more specific, we directly apply gradient updates of the losses to the entries of the IC(0) factors. We initialize training with IC(0) and update only the entries in its prescribed sparsity pattern. In practice, starting from IC(0) and using a small learning rate preserved a positive diagonal and yielded a usable preconditioner throughout.

The gradient of the loss function with respect to the factor entries $L_{ij}$ is computed using automatic differentiation. The update rule follows:
\begin{equation*}
L_{ij}^{(k+1)} = L_{ij}^{(k)} - \alpha \frac{\partial \mathcal{L}}{\partial L_{ij}}
\end{equation*}
where $\mathcal{L}$ is either loss~\eqref{eq:unweighted-frobenius} or loss~\eqref{eq:weighted-frobenius}, $\alpha$ is the learning rate and $k$ denotes the iteration number. Both losses are computed using Hutchinson's stochastic trace estimator~[2] as described in~[5].

We consider $4096\times 4096$ SPD matrix from a cell-centered second-order finite-volume discretization of
\begin{equation*}
-\nabla \cdot \big(k(x, y)\nabla u(x,y)\big) = f(x,y)
\end{equation*}
on a unit square with Dirichlet boundary conditions. The coefficient field $k(x,y)$ is the exponential of a Gaussian random field (global contrast $\approx10$). We draw
$10^3$ random right-hand sides from $\mathcal{N}(0,I)$. Other hyperparameters are: learning rate $10^{-3}$, batch of size $512$, and $10^4$ epochs.

We evaluate three characteristics: (i) the condition number $\kappa(P^{-1}A)$; (ii) CG iteration counts to a fixed relative residual tolerance; and (iii) eigenvalue histograms of $P^{-1}A$. These characteristics consistently indicate that optimizing IC(0) with the weighted objective~\eqref{eq:weighted-frobenius} shifts error away from small-eigenvalue modes, improving conditioning and CG convergence.

\begin{figure*}[ht]
    \vskip 0.2in
    \begin{center}
    \centerline{\includegraphics[width=1.\textwidth]{figures/eigenvalues_distribution.png}}
    \caption{Distributions of eigenvalues of the different matrices. From left to right: initial matrix $A$; preconditioned matrix $P^{-1}A$ with weighted objective~\eqref{eq:weighted-frobenius}; with IC(0); and with unweighted objective~\eqref{eq:unweighted-frobenius}.}
    \label{fig:eigenvalues_distribution}
    \end{center}
    \vskip -0.2in
\end{figure*}

\section{Discussion}

As shown below, weighting the loss by $A^{-1}$ yields a better conditioned $P^{-1}A$ than both IC(0) and the unweighted objective. Table~\ref{tab:condition_number} illustrates the condition number of the preconditioned system with different matrices. \textit{$A$} - initial system matrix. \textit{Weighted Frobenius} - initial matrix preconditioned with factors trained with weighted Frobenius objective~\eqref{eq:weighted-frobenius}; \textit{IC(0)} - initial matrix preconditioned with factors from IC(0); \textit{Frobenius} - initial matrix preconditioned with factors trained with unweighted Frobenius objective~\eqref{eq:unweighted-frobenius}. Numerically, $\kappa(P^{-1}A)$ drops from $11802$ (no preconditioner) to $232$ with IC(0), to $302$ with the unweighted objective, and to $130$ with the weighted objective, i.e., a $\sim1.8\times$ improvement over IC(0) and $\sim2.3\times$ over the unweighted variant.

\begin{table}[ht]
    \caption{Condition number $\kappa(P^{-1}A)$ for different preconditioners: IC(0) baseline, and factors optimized with the unweighted (1) vs. weighted (2) Frobenius objectives.}
    \label{tab:condition_number}
    \centering
    \begin{tabular}{cccc} 
        \toprule
        $A$ & Frobenius & Weighted Frobenius & IC(0) \\
        \midrule
        $11802$ & $302$ & $130$ & $232$ \\
        \bottomrule
    \end{tabular}
\end{table}

The Fig.~\ref{fig:eigenvalues_distribution} illustrates the eigenvalue distributions of the different matrices. Fig.~\ref{fig:cg_iterations} demonstrates CG method convergence with different preconditioners. Fig.~\ref{fig:loss_history} shows the normalized loss convergence during training. The weighted objective yields a tighter eigenvalue cluster for $P^{-1}A$ with fewer extreme small eigenvalues, explaining the lower condition number in Table~\ref{tab:condition_number}.

\begin{figure}[ht]
\vskip 0.2in
\begin{center}
    \centerline{\includegraphics[width=.8\columnwidth]{figures/cg_iterations.png}}
    \caption{CG convergence on a single linear system. Relative residual vs. iterations; same right-hand side for each method. \textit{Weighted Frobenius:} preconditioner optimized with~\eqref{eq:weighted-frobenius}; \textit{IC(0):} baseline incomplete Cholesky; \textit{Frobenius:} preconditioner optimized with~\eqref{eq:unweighted-frobenius}.}
    \label{fig:cg_iterations}
\end{center}
\vskip -0.2in
\end{figure}

\begin{figure}[ht]
    \vskip 0.2in
    \begin{center}
    \centerline{\includegraphics[width=.8\columnwidth]{figures/loss_history.png}}
    \caption{Normalized training loss vs. update steps. Average over stochastic trace batches. \textit{Weighted Frobenius:} objective~\eqref{eq:weighted-frobenius}; \textit{Frobenius:} objective~\eqref{eq:unweighted-frobenius}.}
    \label{fig:loss_history}
    \end{center}
    \vskip -0.2in
\end{figure}

Theoretical analysis shows that the weighted Frobenius functional assigns larger penalties to error components aligned with small eigenvalues of $A$. In applications where $A$ arises from discretizations of elliptic PDEs, these directions often correspond to smooth or low-frequency modes. Thus, minimizing $\|(P - A)A^{-1}\|_F^2$ enforces stronger suppression of such components compared to the plain Frobenius functional. This effect is independent of how the approximation $P$ is parameterized. 
%We emphasize, however, that in unconstrained settings both functionals are trivially minimized by $P = A$, and additional considerations are needed to understand their impact on preconditioning and solver performance.

It is important to note the scope of these observations. The algebraic decomposition is method-agnostic: it describes how the two functionals weight different spectral components of the error, independently of the solver that will use the approximation. The implications for convergence depend on the choice of iterative method and on how the approximation is used. For CG on SPD systems, sensitivity to small eigenvalues makes the connection to our analysis direct. For GMRES or other Krylov subspace methods on nonsymmetric problems, convergence depends on the singular value distribution, the field of values, or the pseudospectrum of the preconditioned matrix. In such settings, suppressing components aligned with small singular values can be beneficial, but it does not constitute a universal guarantee of improved convergence. 

In summary, the weighted Frobenius functional provides a clear, method-independent description of how spectral components of the error are emphasized in structured approximations. Overall, a spectrally weighted Frobenius objective offers a method-agnostic lever to emphasize the right error modes. Its solver-level benefit is clearest for CG on SPD systems and remains problem-dependent elsewhere.



% \section{Conclusion and limitations}

% We have explored a parameter-free approach for improving incomplete Cholesky preconditioners through direct gradient-based optimization of the sparse factors. This method eliminates the need for neural network parameters while potentially achieving improved spectral properties compared to classical IC(0) preconditioners.

% Despite its transparency, the parameter-free approach also faces several limitations. First, the optimization problem is generally nonconvex, so gradient-based methods can only guarantee convergence to stationary points, which may not correspond to globally optimal preconditioners. Second, the number of free variables grows with the number of nonzero entries, which can make the approach computationally demanding and memory-intensive for large-scale problems.

% Primal further work lies in developing theoretical foundations for convergence properties.

% !!!!!!!!!!!!!!!!!!!!!!!!!!!!!!!!!!!!!!!!!!!!!!!!!!!!!!!!!!!!!!!!!!!!!!!!!!!!!!!!!!!!!!!!!!!!!!

% \section{Introduction}
% This document is a model and instructions for \LaTeX.
% Please observe the conference page limits. For more information about how to become an IEEE Conference author or how to write your paper, please visit   IEEE Conference Author Center website: https://conferences.ieeeauthorcenter.ieee.org/.

% \subsection{Maintaining the Integrity of the Specifications}

% The IEEEtran class file is used to format your paper and style the text. All margins, 
% column widths, line spaces, and text fonts are prescribed; please do not 
% alter them. You may note peculiarities. For example, the head margin
% measures proportionately more than is customary. This measurement 
% and others are deliberate, using specifications that anticipate your paper 
% as one part of the entire proceedings, and not as an independent document. 
% Please do not revise any of the current designations.

% \section{Prepare Your Paper Before Styling}
% Before you begin to format your paper, first write and save the content as a 
% separate text file. Complete all content and organizational editing before 
% formatting. Please note sections \ref{AA} to \ref{FAT} below for more information on 
% proofreading, spelling and grammar.

% Keep your text and graphic files separate until after the text has been 
% formatted and styled. Do not number text heads---{\LaTeX} will do that 
% for you.

% \subsection{Abbreviations and Acronyms}\label{AA}
% Define abbreviations and acronyms the first time they are used in the text, 
% even after they have been defined in the abstract. Abbreviations such as 
% IEEE, SI, MKS, CGS, ac, dc, and rms do not have to be defined. Do not use 
% abbreviations in the title or heads unless they are unavoidable.

% \subsection{Units}
% \begin{itemize}
% \item Use either SI (MKS) or CGS as primary units. (SI units are encouraged.) English units may be used as secondary units (in parentheses). An exception would be the use of English units as identifiers in trade, such as ``3.5-inch disk drive''.
% \item Avoid combining SI and CGS units, such as current in amperes and magnetic field in oersteds. This often leads to confusion because equations do not balance dimensionally. If you must use mixed units, clearly state the units for each quantity that you use in an equation.
% \item Do not mix complete spellings and abbreviations of units: ``Wb/m\textsuperscript{2}'' or ``webers per square meter'', not ``webers/m\textsuperscript{2}''. Spell out units when they appear in text: ``. . . a few henries'', not ``. . . a few H''.
% \item Use a zero before decimal points: ``0.25'', not ``.25''. Use ``cm\textsuperscript{3}'', not ``cc''.)
% \end{itemize}

% \subsection{Equations}
% Number equations consecutively. To make your 
% equations more compact, you may use the solidus (~/~), the exp function, or 
% appropriate exponents. Italicize Roman symbols for quantities and variables, 
% but not Greek symbols. Use a long dash rather than a hyphen for a minus 
% sign. Punctuate equations with commas or periods when they are part of a 
% sentence, as in:
% \begin{equation}
% a+b=\gamma\label{eq}
% \end{equation}

% Be sure that the 
% symbols in your equation have been defined before or immediately following 
% the equation. Use ``\eqref{eq}'', not ``Eq.~\eqref{eq}'' or ``equation \eqref{eq}'', except at 
% the beginning of a sentence: ``Equation \eqref{eq} is . . .''

% \subsection{\LaTeX-Specific Advice}

% Please use ``soft'' (e.g., \verb|\eqref{Eq}|) cross references instead
% of ``hard'' references (e.g., \verb|(1)|). That will make it possible
% to combine sections, add equations, or change the order of figures or
% citations without having to go through the file line by line.

% Please don't use the \verb|{eqnarray}| equation environment. Use
% \verb|{align}| or \verb|{IEEEeqnarray}| instead. The \verb|{eqnarray}|
% environment leaves unsightly spaces around relation symbols.

% Please note that the \verb|{subequations}| environment in {\LaTeX}
% will increment the main equation counter even when there are no
% equation numbers displayed. If you forget that, you might write an
% article in which the equation numbers skip from (17) to (20), causing
% the copy editors to wonder if you've discovered a new method of
% counting.

% {\BibTeX} does not work by magic. It doesn't get the bibliographic
% data from thin air but from .bib files. If you use {\BibTeX} to produce a
% bibliography you must send the .bib files. 

% {\LaTeX} can't read your mind. If you assign the same label to a
% subsubsection and a table, you might find that Table I has been cross
% referenced as Table IV-B3. 

% {\LaTeX} does not have precognitive abilities. If you put a
% \verb|\label| command before the command that updates the counter it's
% supposed to be using, the label will pick up the last counter to be
% cross referenced instead. In particular, a \verb|\label| command
% should not go before the caption of a figure or a table.

% Do not use \verb|\nonumber| inside the \verb|{array}| environment. It
% will not stop equation numbers inside \verb|{array}| (there won't be
% any anyway) and it might stop a wanted equation number in the
% surrounding equation.

% \subsection{Some Common Mistakes}\label{SCM}
% \begin{itemize}
% \item The word ``data'' is plural, not singular.
% \item The subscript for the permeability of vacuum $\mu_{0}$, and other common scientific constants, is zero with subscript formatting, not a lowercase letter ``o''.
% \item In American English, commas, semicolons, periods, question and exclamation marks are located within quotation marks only when a complete thought or name is cited, such as a title or full quotation. When quotation marks are used, instead of a bold or italic typeface, to highlight a word or phrase, punctuation should appear outside of the quotation marks. A parenthetical phrase or statement at the end of a sentence is punctuated outside of the closing parenthesis (like this). (A parenthetical sentence is punctuated within the parentheses.)
% \item A graph within a graph is an ``inset'', not an ``insert''. The word alternatively is preferred to the word ``alternately'' (unless you really mean something that alternates).
% \item Do not use the word ``essentially'' to mean ``approximately'' or ``effectively''.
% \item In your paper title, if the words ``that uses'' can accurately replace the word ``using'', capitalize the ``u''; if not, keep using lower-cased.
% \item Be aware of the different meanings of the homophones ``affect'' and ``effect'', ``complement'' and ``compliment'', ``discreet'' and ``discrete'', ``principal'' and ``principle''.
% \item Do not confuse ``imply'' and ``infer''.
% \item The prefix ``non'' is not a word; it should be joined to the word it modifies, usually without a hyphen.
% \item There is no period after the ``et'' in the Latin abbreviation ``et al.''.
% \item The abbreviation ``i.e.'' means ``that is'', and the abbreviation ``e.g.'' means ``for example''.
% \end{itemize}
% An excellent style manual for science writers is \cite{b7}.

% \subsection{Authors and Affiliations}\label{AAA}
% \textbf{The class file is designed for, but not limited to, six authors.} A 
% minimum of one author is required for all conference articles. Author names 
% should be listed starting from left to right and then moving down to the 
% next line. This is the author sequence that will be used in future citations 
% and by indexing services. Names should not be listed in columns nor group by 
% affiliation. Please keep your affiliations as succinct as possible (for 
% example, do not differentiate among departments of the same organization).

% \subsection{Identify the Headings}\label{ITH}
% Headings, or heads, are organizational devices that guide the reader through 
% your paper. There are two types: component heads and text heads.

% Component heads identify the different components of your paper and are not 
% topically subordinate to each other. Examples include Acknowledgments and 
% References and, for these, the correct style to use is ``Heading 5''. Use 
% ``figure caption'' for your Figure captions, and ``table head'' for your 
% table title. Run-in heads, such as ``Abstract'', will require you to apply a 
% style (in this case, italic) in addition to the style provided by the drop 
% down menu to differentiate the head from the text.

% Text heads organize the topics on a relational, hierarchical basis. For 
% example, the paper title is the primary text head because all subsequent 
% material relates and elaborates on this one topic. If there are two or more 
% sub-topics, the next level head (uppercase Roman numerals) should be used 
% and, conversely, if there are not at least two sub-topics, then no subheads 
% should be introduced.

% \subsection{Figures and Tables}\label{FAT}
% \paragraph{Positioning Figures and Tables} Place figures and tables at the top and 
% bottom of columns. Avoid placing them in the middle of columns. Large 
% figures and tables may span across both columns. Figure captions should be 
% below the figures; table heads should appear above the tables. Insert 
% figures and tables after they are cited in the text. Use the abbreviation 
% ``Fig.~\ref{fig}'', even at the beginning of a sentence.

% \begin{table}[htbp]
% \caption{Table Type Styles}
% \begin{center}
% \begin{tabular}{|c|c|c|c|}
% \hline
% \textbf{Table}&\multicolumn{3}{|c|}{\textbf{Table Column Head}} \\
% \cline{2-4} 
% \textbf{Head} & \textbf{\textit{Table column subhead}}& \textbf{\textit{Subhead}}& \textbf{\textit{Subhead}} \\
% \hline
% copy& More table copy$^{\mathrm{a}}$& &  \\
% \hline
% \multicolumn{4}{l}{$^{\mathrm{a}}$Sample of a Table footnote.}
% \end{tabular}
% \label{tab1}
% \end{center}
% \end{table}

% \begin{figure}[htbp]
% \centerline{\includegraphics{fig1.png}}
% \caption{Example of a figure caption.}
% \label{fig}
% \end{figure}

% Figure Labels: Use 8 point Times New Roman for Figure labels. Use words 
% rather than symbols or abbreviations when writing Figure axis labels to 
% avoid confusing the reader. As an example, write the quantity 
% ``Magnetization'', or ``Magnetization, M'', not just ``M''. If including 
% units in the label, present them within parentheses. Do not label axes only 
% with units. In the example, write ``Magnetization (A/m)'' or ``Magnetization 
% \{A[m(1)]\}'', not just ``A/m''. Do not label axes with a ratio of 
% quantities and units. For example, write ``Temperature (K)'', not 
% ``Temperature/K''.

% \section*{Acknowledgment}

% The preferred spelling of the word ``acknowledgment'' in America is without 
% an ``e'' after the ``g''. Avoid the stilted expression ``one of us (R. B. 
% G.) thanks $\ldots$''. Instead, try ``R. B. G. thanks$\ldots$''. Put sponsor 
% acknowledgments in the unnumbered footnote on the first page.

% \section*{References}

% Please number citations consecutively within brackets \cite{b1}. The 
% sentence punctuation follows the bracket \cite{b2}. Refer simply to the reference 
% number, as in \cite{b3}---do not use ``Ref. \cite{b3}'' or ``reference \cite{b3}'' except at 
% the beginning of a sentence: ``Reference \cite{b3} was the first $\ldots$''

% Number footnotes separately in superscripts. Place the actual footnote at 
% the bottom of the column in which it was cited. Do not put footnotes in the 
% abstract or reference list. Use letters for table footnotes.

% Unless there are six authors or more give all authors' names; do not use 
% ``et al.''. Papers that have not been published, even if they have been 
% submitted for publication, should be cited as ``unpublished'' \cite{b4}. Papers 
% that have been accepted for publication should be cited as ``in press'' \cite{b5}. 
% Capitalize only the first word in a paper title, except for proper nouns and 
% element symbols.

% For papers published in translation journals, please give the English 
% citation first, followed by the original foreign-language citation \cite{b6}.

\begin{thebibliography}{00}
\bibitem{b1} Y. Saad, ``Iterative methods for sparse linear systems,'' SIAM, 2003.
\bibitem{b2} M. F. Hutchinson, ``A stochastic estimator of the trace of the influence matrix for Laplacian smoothing splines,'' Communications in Statistics-Simulation and Computation, vol. 18, no. 3, pp. 1059--1076, 1989.
\bibitem{b3} Y. Li, P. Y. Chen, T. Du, and W. Matusik, ``Learning preconditioners for conjugate gradient PDE solvers,'' in International Conference on Machine Learning, 2023, pp. 19425--19439.
\bibitem{b4} V. Trifonov, E. Muravleva, and I. Oseledets, ``Message-passing GNNs fail to approximate sparse triangular factorizations,'' unpublished.
\bibitem{b5} V. Trifonov, A. Rudikov, O. Iliev, Y. M. Laevsky, I. Oseledets, and E. Muravleva, ``Learning from linear algebra: a graph neural network approach to preconditioner design for conjugate gradient solvers,'' unpublished.
\end{thebibliography}

% \vspace{12pt}
% \color{red}
% IEEE conference templates contain guidance text for composing and formatting conference papers. Please ensure that all template text is removed from your conference paper prior to submission to the conference. Failure to remove the template text from your paper may result in your paper not being published.

\end{document}
